\documentclass[../../main.tex]{subfiles}
% NOTE: if your master document isn't called main.tex / isn't one directory up,
% adjust the path in the line above accordingly.

\usepackage{booktabs}              % for \toprule/\midrule/\bottomrule in the
                                    % coefficient tables below (harmless if the
                                    % master preamble already loads it too)
\providecommand{\ad}{\mathrm{ad}}  % falls back to a local definition if the
                                    % master preamble doesn't already provide it

\begin{document}

\section{Cavity Light-Matter Entanglement through Quantum Fluctuations - Passetti et al.}
\label{sec:passetti2022}

\subsection*{Source}
G.\ Passetti, C.\,J.\ Eckhardt, M.\,A.\ Sentef, D.\,M.\ Kennes,
\emph{Cavity Light-Matter Entanglement through Quantum Fluctuations},
arXiv:2212.03011 (2022). 



\subsection{Derivation of the Hamiltonian}
\label{subsec:passetti2022-model}

Within this section, we want to show where the Hamiltonian in equation (1) of the Passetti paper is coming from. The Hamiltonian is
\begin{align}
H&=-t_h\sum_{j=1}^{L}\left(
e^{i\frac{g}{\sqrt L}(a^\dagger+a)}c_j^\dagger c_{j+1}
+\mathrm{H.c.}\right)+\Omega a^\dagger a+H_{\mathrm{int}},
\label{eq:passetti2022-claim}\\
H_{\mathrm{int}}&=U\sum_{j=1}^{L}
\left(n_j-\frac12\right)\left(n_{j+1}-\frac12\right).
\end{align}

\subsection*{Hubbard model}

The book by Nolting -- part seven -- shows a nice derivation of the Hubbard model, which yields (see (2.117) in Nolting part 7)
\begin{align}
H&=\sum_{ij\sigma}T_{ij}c_{i\sigma}^\dagger c_{j\sigma}
%+\frac12 U\sum_{i,\sigma}\hat n_{i\sigma}\hat n_{i,-\sigma}
+ H_{ee}^{(1)}
\notag\\
&=
\sum_{\langle ij\rangle\sigma}T_{ij}c_{i\sigma}^\dagger c_{j\sigma}
%+\frac12 U\sum_{i,\sigma}\hat n_{i\sigma}\hat n_{i,-\sigma},
+ H_{ee}^{(1)}
\end{align}
where we restricted the sum to nearest neighbor hoppings in the second equality. The electronic interaction $H_{ee}^{(1)}$ contains only on site interactions, which means that the overlap integrals
\begin{equation}
\begin{aligned}
V(ij, k \ell)&=\frac{e^2}{4\pi\varepsilon_0}
\int d^3r_1\,d^3r_2\;
w_\sigma^*(r_1-R_i)w_{\sigma'}^*(r_2-R_j)
\frac{1}{|r_1-r_2|}\\
&\hspace{35mm}\times
w_{\sigma'}(r_2-R_k)w_\sigma(r_1-R_\ell).
\end{aligned}
\qquad\begin{gathered}
\sim(2.115)\\
\text{but using}\\
\text{Wannier fcts.}
\end{gathered}
\end{equation}
were simplified to
\begin{equation}
V(ij,k \ell)=
\begin{cases}
U_1&\text{if }i=j=k=\ell,\\
0&\text{else}
\end{cases}
\quad =\delta_{ij}\delta_{jk}\delta_{k\ell}\,U_1.
\end{equation}
Plugging it into the interaction part yields
\begin{align}
H_{ee}^{(1)}
&\overset{\substack{(2.114)\\\text{Nolting}}}{=}
\frac12\sum_{\substack{ijk\ell\\\sigma,\sigma'}}
V^{(1)}(ij,k \ell)
c_{i\sigma}^\dagger c_{j\sigma'}^\dagger
c_{k\sigma'}c_{\ell\sigma}
\notag\\
&=\frac12 U_1\sum_{i,\sigma,\sigma'}
c_{i\sigma}^\dagger c_{i\sigma'}^\dagger
\underbrace{c_{i\sigma'} c_{i\sigma}}_{
\substack{\text{only non-zero for }\sigma'\ne\sigma}}
\notag\\
&=\frac12 U_1\sum_{i,\sigma}
c_{i\sigma}^\dagger c_{i\bar\sigma}^\dagger
c_{i\bar\sigma}c_{i\sigma}
\qquad(\bar\sigma=-\sigma)
\notag\\
&=\frac12 U_1\sum_{i,\sigma}
c_{i\sigma}^\dagger c_{i\sigma}
c_{i\bar\sigma}^\dagger c_{i\bar\sigma}
\notag\\
&=\frac12 U_1\sum_{i,\sigma}n_{i\sigma}n_{i\bar\sigma}.
\label{eq:passetti2022-onsite}
\end{align}

\subsection*{Extended Hubbard model}
Since we will be interested in spinless fermions, we also examine interactions between neighboring sites. This is therefore called \emph{extended} Hubbard model. We only look at processes that do not change the spin of the particles. Thus,
\begin{align}
V(ij,k \ell)&=
\int d^3r_1\,d^3r_2\;
w_\sigma^*(r_1-R_i)w_{\sigma'}^*(r_2-R_j)
\frac{1}{|r_1-r_2|}
\notag\\
&\hspace{35mm}\times
w_{\sigma'}(r_2-R_k) w_\sigma(r_1-R_\ell)
\notag\\
&\underset{\substack{i=\ell,\;j=k\\j=\mathrm{NN}(i)}}{=}
\int d^3r_1\,d^3r_2\;
|w_\sigma(r_1-R_i)|^2|w_{\sigma'}(r_2-R_j)|^2
\frac{1}{|r_1-r_2|}
\notag\\
&\equiv\delta_{i\ell}\delta_{jk}\delta_{i,\mathrm{NN}(j)}\,U_2.
\label{eq:passetti2022-nn-matrix}
\end{align}
This yields another interaction term
\begin{align}
H_{ee}^{(2)}
&\overset{\substack{(2.114)\\\text{Nolting}}}{=}
\frac12\sum_{\substack{ijk\ell\\\sigma,\sigma'}}
V^{(2)}(ij,\ell k)
c_{i\sigma}^\dagger c_{j\sigma'}^\dagger
c_{k\sigma'}c_{\ell\sigma}
\notag\\
&=\frac{U_2}{2}\sum_{\substack{ijk\ell\\\sigma,\sigma'}}
\delta_{i\ell}\delta_{jk}\delta_{i,\mathrm{NN}(j)}
c_{i\sigma}^\dagger c_{j\sigma'}^\dagger
c_{k\sigma'}c_{\ell\sigma}
\notag\\
&=\frac{U_2}{2}\sum_{\langle ij\rangle,\sigma\sigma'}
c_{i\sigma}^\dagger c_{j\sigma'}^\dagger
c_{j\sigma'}c_{i\sigma}
\notag\\
&\underset{i\ne j}{=}\frac{U_2}{2}
\sum_{\langle ij\rangle,\sigma\sigma'}n_{i\sigma}n_{j\sigma'}.
\label{eq:passetti2022-nearest-neighbour}
\end{align}
This finally results in
\begin{equation}
H=
\sum_{\langle ij\rangle\sigma}T_{ij}c_{i\sigma}^\dagger c_{j\sigma}
+\frac12 U_1\sum_{i,\sigma}n_{i\sigma}n_{i\bar\sigma} 
+\frac12 U_2\sum_{\substack{\langle ij\rangle\\\sigma\sigma'}}
n_{i\sigma}n_{j\sigma'}.
\end{equation}

\subsection*{Spinless fermions}

The $n_{i\sigma}n_{i\bar\sigma}$-term vanishes and therefore, with $U_2=U$,
\begin{equation}
H=
\sum_{\langle ij\rangle}T_{ij}c_i^\dagger c_j
+\frac12 U\sum_{\langle ij\rangle}n_i n_j.
\label{eq:passetti2022-spinless}
\end{equation}

\subsection*{Rewrite interaction term}
For periodic boundary conditions and if the number of particles is conserved, we can do the following rewriting
\begin{align}
H_{\text{int,centered}} &= \sum_{\langle ij\rangle}\left(n_i-\frac12\right)\left(n_j-\frac12\right) \\
&=\sum_{\vec r_i} \sum_{\alpha=x,y,z}
\left[n_{\vec r_i} n_{\vec r_i + \hat e_\alpha}
-\frac12(n_{\vec r_i} + n_{\vec r_i + \hat e_\alpha})+\frac14\right]
\notag\\
&=H_{\text{int,uncent.}}
- \frac12 N +\frac B4 - \frac12 \sum_{\vec r_i} (n_{\vec r_i + \hat e_x} + n_{\vec r_i + \hat e_y} +n_{\vec r_i + \hat e_z} ) 
\notag\\
&=H_{\text{int,uncent.}}
- \frac12 N +\frac B4 - \frac32 N =H_{\text{int,uncent.}}
- 2 N +\frac B4,
\notag\\
\label{eq:passetti2022-centering}
\end{align}
where $B$ is the number of bonds of the underlying lattice. Due to particle number conservation, the centered and not centered version of the interaction term only differ by a constant offset. If the particle number is not conserved, this leads to the notion of a chemical potential. This constant shift is irrelevant for the physics of the system. This does not in general hold for open boundary conditions, which can be easily seen for a one dimensional system for instance. The extended Hubbard model for spinless fermions thus reads
\begin{equation}
H=
\sum_{\langle ij\rangle}T_{ij}c_i^\dagger c_j
+\sum_{\langle ij\rangle}\left(n_i-\frac12\right)\left(n_j-\frac12\right)
\quad(+C) .
\label{eq:passetti2022-electronic-centered}
\end{equation}

\subsection*{Introduce light-matter interaction via scheme presented by Dmytruk and Schiro}
The general idea to introduce the coupling between light and matter that Schiro and Dmytruk implemented goes as follows. The light-matter interaction is introduced via minimal coupling $-i\nabla \rightarrow -i\nabla e \bm A(\bm r)$ in the continuum, which yields
\begin{align}
    \mathcal H_C &= \mathcal H_0 + \mathcal H_{ee} + \mathcal H_{\text{ph}} \qquad \text{with} \\
    \mathcal H_0 &= \int d \bm r \, \psi^\dagger(\bm r) \left( \frac{(-i\nabla + e\bm A(\bm r))^2}{2m} + V(\bm r) \right) \psi(\bm r) , \\
    \mathcal H_{ee} &= \int d \bm r d \bm r \, \psi^\dagger(\bm r) \psi^\dagger(\bm r') U(\bm r - \bm r') \psi(\bm r') \psi(\bm r) , \\
    \mathcal H_{\text{ph}} &= \int d \bm r \, \bm \Pi(\bm r)^2 + (\nabla \times \bm A(\bm r))^2 .
\end{align}
They show, that this replacement can also be achieved by performing a unitary transformation on the matter degrees of freedom of the system only, i.e., 
\begin{align}
    \mathcal H_C &= \mathcal H_{\text{ph}} + \mathcal U^\dagger (\mathcal H_0 + \mathcal H_{ee}) \mathcal U \qquad \text{where} \\
    \mathcal U(\chi) &= \exp\left[ i(a+a^\dagger) \int d\vec r \psi^\dagger(\bm r) \chi(\bm r)\psi(\bm r) \right], \\
    \nabla \chi(\bm r) &= e \bm A_0(\bm r) .
\end{align} 
They project the matter Hamiltonian to a set of low energy states and then apply the unitary transformation $\mathcal U \rightarrow U$, which itself has also been projected. 

There are apparently drawbacks with the approach presented by Dmytruk and Schiro but they do not manifest for 1d systems with only one band. For a complemented approach, look at the paper \emph{Electromagnetic coupling in tight-binding models for strongly correlated light and matter} by Li and Eckstein et al.

\paragraph{Unitary transformation on the electronic part of the Hamiltonian in the Coulomb gauge}

\begin{equation}
U^\dagger(\chi)c_{\bm R\mu}U(\chi)
=e^{i(a+a^\dagger)\chi(\bm R)}c_{\bm R\mu}
\qquad\text{(S29, 36--38)}
\label{eq:passetti2022-local-phase}
\end{equation}
\begin{equation}
\Rightarrow\quad U^\dagger c_{\bm R\mu}^\dagger c_{\bm R'\mu'}U
=e^{i(a+a^\dagger)(\chi(\bm R')-\chi(\bm R))}
c_{\bm R\mu}^\dagger c_{\bm R'\mu'}
\end{equation}
where
\begin{align}
\chi(\bm R')-\chi(\bm R)
&\overset{\text{(S14)}}{=}
\int_{\gamma(\bm r^*,\bm R')}A(\bm s)\,d\bm s
-\int_{\gamma(\bm r^*,\bm R)}A(\bm s)\,d\bm s
\notag\\
&=\int_{\bm R}^{\bm R'}A(\bm s)\,d\bm s
+\int_{\Gamma(\bm r^*,\bm R,\bm R')}A(\bm s)\,d\bm s.
\label{eq:passetti2022-path-difference}
\end{align}

% \begin{center}
% \setlength{\unitlength}{1pt}
% \begin{picture}(260,95)
% \put(20,17){\circle*{3}}
% \put(130,75){\circle*{3}}
% \put(235,17){\circle*{3}}
% \put(9,3){$\bm R$}
% \put(120,81){$\bm R'$}
% \put(238,3){$\bm r^{\,*}$}
% \qbezier(20,17)(80,55)(130,75)
% \qbezier(130,75)(190,90)(235,17)
% \qbezier(20,17)(45,43)(78,18)
% \qbezier(78,18)(130,-4)(170,28)
% \qbezier(170,28)(218,53)(235,17)
% \put(110,40){$\Gamma(\bm r^*,\bm R,\bm R')$}
% \put(125,0){$\gamma(\bm r^*,\bm R)$}
% \put(180,76){$\gamma(\bm r^*,\bm R')$}
% \end{picture}
% \end{center}

In one dimension, the line integral simplifies to 

\begin{align}
\Rightarrow\quad\chi(R')-\chi(R)
&=\int_R^{R'}A(s)\,ds
\notag\\
&\underset{\substack{A(s)=A\;\text{dipole approx.}\\+\,1d\;\mathrm{NN}}}{=}
A(j+1-j)=A
\end{align}
\begin{equation}
\Rightarrow\quad U^\dagger c_j^\dagger c_{j+1}U
=e^{ig(a+a^\dagger)}c_j^\dagger c_{j+1}
\qquad\sim\text{(S44)}
\end{equation}
and
\begin{equation}
U^\dagger n_j U=n_j,\qquad\text{since }\int_R^R A(s)\,ds=0
\end{equation}
\begin{equation}
\begin{aligned}
\Rightarrow\quad H_C
&\overset{\substack{\text{S44 in 1d}}}{=}
\omega_c a^\dagger a
+\sum_j t_{j,j+1}\left(
e^{ig(a+a^\dagger)}c_j^\dagger c_{j+1}+\mathrm{H.c.}\right)
\end{aligned}
\label{eq:passetti2022-coulomb}
\end{equation}

Combining \cref{eq:passetti2022-coulomb,eq:passetti2022-electronic-centered} finally yields
\begin{equation}
\begin{aligned}
H=&-t_h\sum_{j=1}^{L}\left(
e^{i\frac{g}{\sqrt L}(a^\dagger+a)}c_j^\dagger c_{j+1}
+\mathrm{H.c.}\right)+\Omega a^\dagger a &\qquad\text{(P1)}\\
&+U\sum_{j=1}^{L}\left(n_j-\frac12\right)
\left(n_{j+1}-\frac12\right)
&\qquad\text{(P8)}
\end{aligned}
\label{eq:passetti2022-derived-model}
\end{equation}

%%%%%%%%%%%%%%%%%%%%%%%%%%%%%%%%%%%%%%%%%%%%%%%%%%%%%%%%%%%%%%%%%%%%%%%%%

\subsection{Appendix I: High-frequency expansion} 
\label{subsec:passetti2022-high-frequency}

\subsection*{Roadmap}
This note reconstructs, with every intermediate step shown, the road from the
main-text Hamiltonian Eq.~(1) to the Supplemental Material's transformed
Hamiltonian Eq.~(S3) and renormalized frequency Eq.~(S4):
\begin{center}
Eq.~(1) $\xrightarrow{\text{Taylor expand Peierls phase, }O(g^2)}$ Eq.~(S1)
$\xrightarrow{\text{displace}} h'(j,t)
\xrightarrow{\text{squeeze}}$ Eq.~(S3), Eq.~(S4).
\end{center}
The Hadamard-lemma machinery used in Steps~3--4 is general-purpose and is
derived in full in the note \emph{Similarity Transformations of Operators:
Computing $e^{A}Be^{-A}$} in the general reference-notes project; here we only
recall the results needed and apply them to this specific Hamiltonian. Main-text equations are cited as Eq.~(n); equations
from the Supplemental Material are cited as Eq.~(Sn).

\subsection*{Starting point: the model Hamiltonian}

Eq.~(1) of the main text describes a chain of $L$ spinless fermions,
nearest-neighbour hopping $t_h$ and interaction $U$, Peierls-coupled to a
single cavity mode of frequency $\Omega$ via the (spatially uniform, $q_x=0$)
mode with coupling $g$:
\begin{equation}
H=-t_h\sum_{j=1}^{L}\left(e^{i\frac{g}{\sqrt L}(a^\dagger+a)}c_j^\dagger c_{j+1}
+\mathrm{h.c.}\right)
+U\sum_{j=1}^{L}\left(n_j-\tfrac12\right)\left(n_{j+1}-\tfrac12\right)
+\Omega a^\dagger a. \label{eq:maintext1}
\end{equation}
Only the hopping term carries $g$-dependence; the interaction and bare photon
terms are untouched by the Peierls substitution.

\subsection*{Step 1: expanding the Peierls phase to quadratic order in \textit{g}}
\label{sec:step1-expansion}

Write $\theta\equiv\dfrac{g}{\sqrt L}(a^\dagger+a)$. Since $\theta^\dagger=\theta$, $(e^{i\theta})^\dagger
=e^{-i\theta}$, so the hopping term of Eq.~\eqref{eq:maintext1} is
\begin{equation}
-t_h\sum_j\left(e^{i\theta}c_j^\dagger c_{j+1}+e^{-i\theta}c_{j+1}^\dagger c_j\right).
\end{equation}
Taylor-expand each exponential to $O(\theta^2)$:
\begin{equation}
e^{\pm i\theta}=1\pm i\theta-\frac{\theta^2}{2}+O(\theta^3).
\end{equation}
Substituting and collecting by power of $\theta$:
\begin{align}
&-t_h\sum_j\left[\left(1+i\theta-\tfrac{\theta^2}2\right)c_j^\dagger c_{j+1}
+\left(1-i\theta-\tfrac{\theta^2}2\right)c_{j+1}^\dagger c_j\right]\notag\\
=\;&\underbrace{-t_h\sum_j\left(c_j^\dagger c_{j+1}+c_{j+1}^\dagger c_j\right)}_{\theta^0}
\;\underbrace{-\,it_h\theta\sum_j\left(c_j^\dagger c_{j+1}-c_{j+1}^\dagger c_j\right)}_{\theta^1}
\;\underbrace{+\,\frac{t_h\theta^2}{2}\sum_j\left(c_j^\dagger c_{j+1}+c_{j+1}^\dagger c_j\right)}_{\theta^2}
+O(\theta^3). \label{eq:taylorexpand}
\end{align}

Define the kinetic-energy and current operators exactly as in the main text,
\begin{equation}
T\equiv-t_h\sum_j\left(c_j^\dagger c_{j+1}+c_{j+1}^\dagger c_j\right),
\qquad
J\equiv-it_h\sum_j\left(c_j^\dagger c_{j+1}-c_{j+1}^\dagger c_j\right).
\end{equation}
Using $c_j^\dagger = \frac{1}{\sqrt{L}}\sum_{k} e^{-ikj} c_k^\dagger$ with $k=\frac{2\pi m}{L}, m=0,\dots L-1$:
\begin{equation}
T=-2t_h\sum_k \cos(k) c_k^\dagger c_k,
\qquad
J=2t_h\sum_k \sin(k) c_k^\dagger c_k.
\end{equation}
The dispersion relation for the non-interacting case ($U=g=0$) is $\epsilon(k) = -2t_h \cos(k)$ and thus the group velocity of each electron is (see Bartelmann for example) $v_g(k) = \frac{\partial \epsilon(k)}{\partial k} = 2t_h \sin(k)$. Therefore, $J$ is truly a current operator (remember that $e=1$). Alternatively, this identification can be argued since $J$ couples linearly to the vector potential $A\propto a^\dagger + a$ in the Hamiltonian, just as in classical electrodynamics (compare to Bartelmann part 2).
The $\theta^0$ term is exactly $T$; the $\theta^2$ term equals
$-\tfrac{\theta^2}{2}T$, since $t_h\sum_j(c_j^\dagger c_{j+1}+\mathrm{h.c.})=-T$.
So Eq.~\eqref{eq:taylorexpand} becomes
\begin{equation}
T+\theta J-\frac{\theta^2}{2}T+O(\theta^3).
\end{equation}
Substituting $\theta=\dfrac{g}{\sqrt L}(a^\dagger+a)$,
\begin{equation}
\theta J=\frac{g}{\sqrt L}(a^\dagger+a)J,\qquad
\frac{\theta^2}{2}T=\frac{g^2}{2L}(a^\dagger+a)^2T,
\end{equation}
and adding back the untouched interaction and photon terms from
Eq.~\eqref{eq:maintext1} gives exactly Eq.~(S1) of the Supplemental Material:
\begin{equation}
\boxed{\,H^{g^2}=\Omega a^\dagger a+T+H_{\mathrm{int}}
+\frac{g}{\sqrt L}(a^\dagger+a)J-\frac{g^2}{2L}(a^\dagger+a)^2T\,}
\qquad\text{[Eq.~(S1)]}
\label{eq:Hg2}
\end{equation}
with $H_{\mathrm{int}}=U\sum_j(n_j-\tfrac12)(n_{j+1}-\tfrac12)$. This
truncation is controlled: the dropped terms are $O(\theta^3)=O(g^3)$, and are
parametrically small in the regime the paper works in, $gt_h\ll\Omega$.

\subsection*{Step 2: reducing to a family of photon-only Hamiltonians}
\label{sec:step2-reduction}

$T$ and $J$ commute, $[T,J]=0$ since both are diagonal in the same momentum-space
single-particle basis, and both trivially commute with the photon operators
$a,a^\dagger$ (different tensor factors). Furthermore, one can show that the interaction term 
$H_{\mathrm{int}}$ transforms under displacement and squeezing as 
$H_{\mathrm{int}} \rightarrow H_{\mathrm{int}} + O(\frac{1}{\Omega})$. 
Since the terms $O(\frac{1}{\Omega})$ are also dropped, $H_{\mathrm{int}}$ transforms trivially. 
So, restrict to a simultaneous
eigenspace $|j,t\rangle$ of $(J,T)$ with eigenvalues $(j,t)$, everything in
Eq.~\eqref{eq:Hg2} except $H_{\mathrm{int}}$ becomes an ordinary
(c-number-parametrized) bosonic Hamiltonian:
\begin{equation}
h(j,t)=\Omega a^\dagger a+\frac{g}{\sqrt L}j\,(a^\dagger+a)
-\frac{g^2}{2L}t\,(a^\dagger+a)^2. \label{eq:hjt}
\end{equation}
Equation~\eqref{eq:hjt} is accordingly the object to diagonalize in the photon
sector only.

\subsection*{Step 3: displacement transformation}
\label{sec:step3-displacement}

Define $\mathcal D\equiv e^{S^d}$, $S^{d}=\alpha(a^\dagger-a)$ with $\alpha$
a real number (later fixed as a function of $j$). Since $[S^d,a]=-\alpha$ is a
c-number, the Hadamard series truncates exactly:
\begin{equation}
e^{S^d}a\,e^{-S^d}=a-\alpha,\qquad e^{S^d}a^\dagger e^{-S^d}=a^\dagger-\alpha.
\end{equation}
Applying this to Eq.~\eqref{eq:hjt} and keeping terms to $O(g^2)$ (dropping the
$O(g^3)$ piece generated by displacing the $(a^\dagger+a)^2$ term, consistent
with the truncation of Eq.~\eqref{eq:Hg2} itself):
\begin{equation}
e^{S^d}h(j,t)e^{-S^d}=\Omega a^\dagger a
-\Omega\alpha(a^\dagger+a)+\frac{g}{\sqrt L}j(a^\dagger+a)
-\frac{g^2}{2L}t(a^\dagger+a)^2
+\left(\Omega\alpha^2-\frac{2g}{\sqrt L}j\alpha\right)+O(g^3).
\end{equation}
Killing the linear term fixes
\begin{equation}
\boxed{\;\alpha(j)=\frac{g}{\Omega\sqrt L}\,j\;}\qquad\text{[cf. Eq.~(S2)]},
\end{equation}
and the constant collapses to $\Omega\alpha^2-2\Omega\alpha^2=-\Omega\alpha^2
=-\dfrac{g^2j^2}{\Omega L}$, giving
\begin{equation}
h'(j,t)\equiv e^{S^d}h(j,t)e^{-S^d}
=\Omega a^\dagger a-\frac{g^2t}{2L}(a^\dagger+a)^2
-\frac{g^2j^2}{\Omega L}+O(g^3). \label{eq:hprime}
\end{equation}

\subsection*{Step 4: squeezing transformation}
\label{sec:step4-squeezing}

Define $\mathcal S\equiv e^{S^{\mathrm{sq}}}$, $S^{\mathrm{sq}}=\tfrac{\xi}{2}
(a^2-a^{\dagger2})$; $\xi=\xi(t)$ is to be fixed by requiring the transformed
Hamiltonian is diagonal in $a^\dagger a$.

\paragraph{Transforming $a$ and $a^\dagger$ (Hadamard lemma).}
Using $[AB,C]=A[B,C]+[A,C]B$: $[S^{\mathrm{sq}},a]=\xi a^\dagger$,
$[S^{\mathrm{sq}},a^\dagger]=\xi a$. Iterating (even orders return $\xi^{2k}a$,
odd orders return $\xi^{2k+1}a^\dagger$; check for instance via Python's Sympy) and summing the resulting series:
\begin{equation}
e^{S^{\mathrm{sq}}}a\,e^{-S^{\mathrm{sq}}}=a\cosh\xi+a^\dagger\sinh\xi,
\qquad
e^{S^{\mathrm{sq}}}a^\dagger e^{-S^{\mathrm{sq}}}=a^\dagger\cosh\xi+a\sinh\xi.
\label{eq:aTsq}
\end{equation}

\paragraph{Transforming $a^\dagger a$ and $(a^\dagger+a)^2$.}
Using $e^S(XY)e^{-S}=(e^SXe^{-S})(e^SYe^{-S})$
on Eq.~\eqref{eq:aTsq} and normal-ordering with $aa^\dagger=a^\dagger a+1$:
\begin{align}
e^{S^{\mathrm{sq}}}(a^\dagger a)e^{-S^{\mathrm{sq}}}
&=a^\dagger a\,\cosh(2\xi)+\sinh^2\xi+\tfrac12\sinh(2\xi)\big(a^2+a^{\dagger2}\big),
\label{eq:numtransform}\\
e^{S^{\mathrm{sq}}}(a^\dagger+a)^2e^{-S^{\mathrm{sq}}}
&=e^{2\xi}\big(a^{\dagger2}+a^2+2a^\dagger a+1\big). \label{eq:sqtransform}
\end{align}
The rewriting was done using $\cosh^2(x) + \sinh^2(x) = \cosh(2x)$, $2\sinh(x)\cosh(x) = \sinh(2x)$ and the definition of both functions via the exponential function.

\paragraph{Assembling and diagonalizing.} Substituting
Eqs.~\eqref{eq:numtransform}--\eqref{eq:sqtransform} into $h'(j,t)$,
Eq.~\eqref{eq:hprime}, and collecting terms:
\begin{center}
\begin{tabular}{ll}
\toprule
Operator & Coefficient \\
\midrule
$a^\dagger a$ & $\Omega\cosh(2\xi)-\dfrac{g^2t}{L}e^{2\xi}$ \\[2mm]
$a^2+a^{\dagger2}$ & $\dfrac{\Omega}{2}\sinh(2\xi)-\dfrac{g^2t}{2L}e^{2\xi}$ \\[2mm]
const. & $\Omega\sinh^2\xi-\dfrac{g^2t}{2L}e^{2\xi}-\dfrac{g^2j^2}{\Omega L}$ \\
\bottomrule
\end{tabular}
\end{center}
Requiring the $(a^2+a^{\dagger2})$ coefficient to vanish,
$\Omega\sinh(2\xi)=\dfrac{g^2t}{L}e^{2\xi}$, solved via
$e^{4\xi}=\left(1-\dfrac{2g^2t}{\Omega L}\right)^{-1}$:
\begin{equation}
\xi(t)=-\frac14\ln\!\left(1-2\frac{g^2}{\Omega}\frac tL\right).
\label{eq:xiresult}
\end{equation}

\paragraph{Remark: discrepancy with the printed Supplemental Material formula.}
The Supplemental Material states $\xi(t)=\tfrac12\ln\!\left(1-2\tfrac{g^2}{\Omega}
\tfrac tL\right)$, a factor of $-2$ different from Eq.~\eqref{eq:xiresult}. This
looks like a transcription slip in the supplement. It does not affect
Eq.~(S4) below, which depends only on $\xi^2$-type combinations.

\subsection*{Assembling the result: Eq.~(S3) and the renormalized frequency Eq.~(S4)}
\label{sec:assembling}

Plug Eq.~\eqref{eq:xiresult} back into the $a^\dagger a$ coefficient. With
$\frac{g^2 t}{L} = \frac{\Omega}{2}(1-e^{-4\xi})$:
\begin{equation}
\tilde\Omega=\Omega\cosh(2\xi)-\frac{g^2t}{L}e^{2\xi}
=\frac{\Omega}{2}\left(e^{2\xi} + e^{-2\xi}\right)
-\frac{g^2 t}{L}e^{2\xi}
=\Omega\,e^{-2\xi},
\end{equation}
i.e.
\begin{equation}
\boxed{\;\tilde\Omega=\Omega\sqrt{1-2\dfrac{g^2}{\Omega}\dfrac{t}{L}}\;}
\qquad\text{[Eq.~(S4), with }t\to\langle T\rangle\text{]}.
\label{eq:OmegaTilde}
\end{equation}
This matches the paper exactly. The displaced and squeezed Hamiltonian in the photon sector thus reads

\begin{equation}
    \tilde \Omega a^\dagger a + \frac{\tilde \Omega - \Omega}{2} - \frac{g^2 j^2}{\Omega L} = \tilde \Omega a^\dagger a + O (\frac{1}{\Omega})
\end{equation}

Reattaching the fermionic sector, the fully assembled transformed
Hamiltonian is Eq.~(S3):
\begin{align}
\tilde H^{g^2}&=\tilde\Omega a^\dagger a+T+H_{\mathrm{int}} \\
&=\tilde\Omega \beta^\dagger\beta+T+H_{\mathrm{int}}
\qquad\text{[Eq.~(S3)]},
\end{align}
where $\beta$ is the new (squeezed, displaced) photon annihilation operator,
and $\tilde\Omega$ depends on the fermionic sector only through $t=\langle T\rangle$
(and reduces to $\Omega$ at $O(g^0)$, as it must). The first line, describes the transformed Hamiltonian in the original photon basis, whereas the second line describes it in the displaced and squeezed basis, and therefore $\beta = e^{S^\mathrm{sq}} e^{S^\mathrm{d}} a e^{-S^\mathrm{sq}} e^{-S^\mathrm{d}}$. As emphasized in
Sec.~\ref{sec:step2-reduction}, the appearance of $H_{\mathrm{int}}$ unchanged
here is a leading-order statement in $1/\Omega$, not an exact operator
identity.

\subsection*{See also}
For the general-purpose machinery reused in Steps~\ref{sec:step3-displacement}
and \ref{sec:step4-squeezing} (the Hadamard lemma, the multiplicativity trick,
the dagger shortcut, and the equivalent differential-equation/generator-flow
formulation), see the note \emph{Similarity Transformations of Operators:
Computing $e^{A}Be^{-A}$} in the general reference-notes project.

\end{document}
